\section{Experimental design}
\label{sec:experiments}

Our central question is whether particle-pair information in Particle
Transformer (ParT)~\cite{Qu:2022mxj} should only control attention weights
or also contribute to the messages passed between particles. We examine
this question primarily on offline JetClass inputs, treating shared versus
layerwise attention bias and the addition of relation-conditioned value
messages as distinct architectural choices. Standard ParT pair features
provide the reference setting; selected additional physics-derived
relations form a secondary extension. An earlier exploratory study on a
synthetic HLT-like reduced-input view informed that feature selection and
provides supporting comparisons. It is not a trigger-performance benchmark
or the primary basis for the offline architectural conclusions.

\subsection{Offline dataset and split roles}
\label{sec:experimental-data}

We use the ten-class JetClass dataset~\cite{Qu:2022mxj,JetClass}. The labels
are QCD, $H\to b\bar b$, $H\to c\bar c$, $H\to gg$, $H\to4q$,
$H\to q\bar q\ell\nu$, $W\to q\bar q'$, $Z\to q\bar q$,
$t\to bq\bar q'$, and $t\to b\ell\nu$. Offline inputs are the supplied
reconstructed constituents. Each configuration uses the same balanced
sample of one million training jets, or 100,000 jets per class. Our
principal evaluation targets the official 20-million-jet test set, with
two million jets per class. This increases evaluation statistics, not
training-set size. Table~\ref{tab:experimental-splits} distinguishes the
development, earlier test, and full-test samples.

\begin{table}[htbp]
\centering
\small
\begin{tabularx}{\linewidth}{@{}lrX@{}}
\toprule
Sample & Jets & Role \\
\midrule
Training & 1,000,000 & Fit model weights and relation normalization. \\
Checkpoint validation & 125,000 & Select checkpoints and determine early stopping. \\
Comparison validation & 125,000 & Compare configurations; used for the earlier reduced-view relation screen. \\
Earlier test & 500,000 & Supporting reduced-view evaluation and source of provisional offline draft values. \\
Official full test & 20,000,000 & Primary offline evaluation of frozen checkpoints (pending). \\
\bottomrule
\end{tabularx}
\caption{Split roles. The first four samples are disjoint, class-balanced
subsets from the original campaigns. The full-test protocol checks
source-file disjointness against them before inference. The official test
does not enlarge the one-million-jet training sample.}
\label{tab:experimental-splits}
\end{table}

Inputs contain at most 128 constituents per jet, with padded entries
masked. Jet axes, relative kinematics, four-vectors, and pair descriptors
are derived from these inputs. Normalization statistics are fitted only
on training data. Every offline model is trained from scratch; no
HLT-trained weights, teacher predictions, or privileged auxiliary inputs
are used.

The earlier reduced-view and offline campaigns share jet identities and
split assignments. Their comparison is therefore across representations
of the same underlying samples, not independent development datasets.
The earlier offline test identities may also have appeared in previous
repository experiments; we do not describe that sample as globally
untouched. The full-test evaluation is a separate assessment of frozen
checkpoints.

\subsection{Common backbone and experimental factors}
\label{sec:experimental-models}

All configurations retain a common ParT-style particle encoder and
class-attention readout~\cite{Qu:2022mxj}. Our implementation supplies 17
particle features describing kinematics, charge, particle identification,
impact parameters, and their uncertainties. The particle embedding has
widths $(128,512,128)$, followed by eight particle-attention blocks with
eight heads and two class-attention blocks. The standard pair encoder uses
hidden widths $(64,64,64)$. The reference pair inputs are the logarithms of
four standard ParT quantities: angular separation, relative transverse-momentum
scale, momentum-sharing fraction, and pair invariant mass squared. The
baseline and its extensions use the same explicit-pair interface and
preprocessing, avoiding an input-path change between comparisons.

Our central architectural factors are \emph{bias sharing} and the
\emph{value-message pathway}. Shared bias reuses one head-specific pair-bias
tensor across particle-attention layers; layerwise bias instead uses
independent final bias projections for each layer. In both cases, the
learned pair stem is evaluated once and reused, not recurrently updated.
Changing bias sharing alters how pair information controls attention
weights. EdgeValue additionally lets that information contribute to the
message being aggregated.

We use \emph{EdgeValue} as our label for a ParT adaptation of established
relation-aware value aggregation. Shaw et al.~\cite{Shaw2018Relative}, Eq.~(3), add an edge
vector inside the attention-weighted value sum; we adopt this value-side
mechanism rather than introduce a new general attention operator.
For head $h$ in layer $\ell$, our parameterization has the form
\begin{align}
\alpha_{ij}^{\ell,h}
&= \operatorname{softmax}_{j}\left(
\frac{(q_i^{\ell,h})^\top k_j^{\ell,h}}{\sqrt{d_h}}
+ b_{ij}^{\ell,h}\right), \\
m_i^{\ell,h}
&= \sum_j \alpha_{ij}^{\ell,h}
\left(v_j^{\ell,h}+W_{\mathrm{edge}}^{\ell,h}r_{ij}\right).
\label{eq:experimental-edgevalue}
\end{align}
Here $r_{ij}$ is the learned pair-stem representation and
$W_{\mathrm{edge}}^{\ell,h}$ is a layer- and head-specific, bias-free linear
projection. The vector $W_{\mathrm{edge}}^{\ell,h}r_{ij}$ corresponds to
the relation-vector term $a_{ij}^{V}$ in
Shaw et al.~\cite{Shaw2018Relative}, parameterized here using learned
particle-pair descriptors. Our score retains a ParT-style scalar pair bias
rather than their query-dependent relative-key term. Without EdgeValue,
the additional value vector is absent. Shared
bias makes $b_{ij}^{\ell,h}$ independent of $\ell$; it does not remove
pairwise bias. The value term uses the same attention weights as the
ordinary value aggregation, while the existing output projection and
class-attention readout are retained. Its linear projection can be applied
after the attention-weighted sum of pair representations, avoiding an
explicit per-edge, per-head value-vector tensor.

\subsection{Offline architectural comparisons}
\label{sec:offline-design}

The standard-feature comparisons in Table~\ref{tab:offline-design} form
a planned $2\times2$ architecture study: shared or layerwise bias, each
with or without EdgeValue. Holding the pair features fixed is intended to
separate the contribution of value messages from a change in bias sharing.
The selected-feature rows then ask whether additional pair descriptors
benefit these architectures. The full seven-row matrix is not a complete
$2\times2\times2$ factorial study: the selected-feature, shared-bias model
without EdgeValue is absent.

The original campaign trained four configurations from scratch with three
matched random seeds: the standard shared-bias baseline; standard features
with layerwise bias and EdgeValue; selected features with layerwise bias;
and selected features with layerwise bias and EdgeValue. All twelve
checkpoints are retained regardless of validation performance. Its
predeclared primary comparison was selected features with layerwise bias
and EdgeValue versus the baseline. That comparison measures a combined
design change, not the isolated effect of EdgeValue.

Within those four configurations, the selected-feature layerwise pair
with versus without EdgeValue holds both the pair descriptors and bias
sharing fixed and directly tests the value pathway. In contrast, the
standard-feature EdgeValue model differs from the baseline in both bias
sharing and value messages. Completing the standard-feature
$2\times2$ comparison requires the two additional standard-feature rows;
the selected shared-bias EdgeValue row extends the bias-sharing comparison
to the augmented feature set. For shared-bias EdgeValue models, the
layer- and head-specific value projections remain unchanged.

\begin{table}[htbp]
\centering
\small
\begin{tabular}{@{}llcl@{}}
\toprule
Pair features & Pair bias & EdgeValue & Training status \\
\midrule
Standard & Shared & No & Completed \\
Standard & Layerwise & No & Planned \\
Standard & Shared & Yes & Planned \\
Standard & Layerwise & Yes & Completed \\
Selected & Layerwise & No & Completed \\
Selected & Shared & Yes & Planned \\
Selected & Layerwise & Yes & Completed \\
\bottomrule
\end{tabular}
\caption{Offline experimental matrix. Each row uses the same three training seeds.
Selected means standard features plus PT, TRACK, and REGION. Completed
refers to training, not completion of the official full-test evaluation.
The three new rows are prospective ablations, with no results yet.}
\label{tab:offline-design}
\end{table}

The three additional ablations were specified after inspection of the
earlier 500,000-jet offline results; they were not part of the original
four-model predeclared study. They are to use the same training protocol
and matched seeds. Their nine checkpoints still require training and
full-test evaluation. The full-test evaluation of the twelve existing
checkpoints is underway; no such checkpoint is retrained or selected using
its full-test results. Until those metrics are available, the results
section uses explicitly labeled 500,000-jet drafting placeholders for the
four trained configurations, and TBD for the untrained rows.

\subsection{Additional pair descriptors and exploratory selection}
\label{sec:selected-relations}

The secondary feature extension retains the four standard ParT inputs and
adds three encoded relation families. PT describes endpoint momentum
fractions, ranks, and directed momentum differences. TRACK describes
impact parameters, their uncertainties and significances, and pairwise
compatibility. REGION describes shared cluster membership, common-ancestor
properties, and endpoint cluster information at exclusive resolutions
$K=2,4,8$. Its deterministic, beam-free angular tree is reconstructed from
the available constituents, not generator ancestry or truth vertices.
The respective encoded widths are 8, 12, and 12, giving
$4+8+12+12=36$ pair-input channels before the shared pair stem.
All descriptors are recomputed from the offline inputs, with normalization
refitted on offline training data.

PT+TRACK+REGION was selected before the offline campaign using comparison
validation in an earlier synthetic HLT-like study. That study screened 21
shared-bias configurations and examined architecture variants and controls
with three matched seeds. We retain it as an exploratory characterization
of pair information under a fixed input degradation, not evidence that
the selected families are universally optimal or required for EdgeValue.
The selection protocol, synthetic-view definition, and complete supporting
results are given in Appendix~\ref{sec:hlt-support}. This preserves the
origin of the feature choice without making the offline architecture study
depend on a claim of realistic trigger simulation.

\subsection{Training and checkpoint selection}
\label{sec:training-protocol}

All completed configurations use unweighted ten-class cross-entropy and
AdamW with learning rate $10^{-3}$, $(\beta_1,\beta_2)=(0.9,0.999)$,
weight decay $10^{-4}$, and gradient-norm clipping at 1.0. Microbatches of
64 jets are accumulated over two steps, giving a nominal effective batch
size of 128. A final incomplete accumulation group is retained and
normalized by its actual event count. Training uses a maximum of 40 epochs,
with a minimum of 12 epochs and early-stopping patience of eight epochs.
Completed offline training used BF16 mixed precision. Offline configurations
use the same three training seeds and, within each matched seed, the same
deterministic epoch-wise ordering of jet identities.

The learning rate increases linearly for the first 5\% of the planned
optimizer updates, rounded upward, then follows a cosine schedule to
$10^{-5}$. With one million training jets, the schedule contains 7,813
updates per epoch, 312,520 planned updates, and 15,626 warm-up updates.
Early stopping does not shorten or rescale this planned schedule.

After each epoch, checkpoints are compared only on checkpoint validation.
Among epochs within $10^{-4}$ absolute accuracy of the best observed
accuracy, the checkpoint with lowest cross-entropy is preferred, with the
earliest epoch breaking exact ties. Patience resets when the current epoch
becomes the preferred checkpoint. Comparison validation is evaluated for
the selected checkpoint and is not used to choose another epoch. This
protocol is shared, but early stopping means that realized training lengths
need not be equal. Added relation encoders and value projections also
change parameter counts and computational cost; the comparisons are not
described as universally capacity- or compute-matched.

\subsection{Evaluation and uncertainty}
\label{sec:evaluation-protocol}

We report classification accuracy, cross-entropy, macro one-versus-rest
AUC, and class-specific QCD rejection. Macro AUC is the unweighted mean
of the ten classwise AUCs, including QCD, computed separately for each
checkpoint. Each classwise AUC compares that class with all remaining
classes using its multiclass softmax probability as the score. This differs
from the signal-versus-QCD discriminant used for background rejection.
Macro AUC is reported on the 0--1 scale; its multi-seed standard deviation
is computed across the three per-seed macro AUCs, not across classes.
These additional reported metrics do not change the original accuracy and
cross-entropy selection rule. Accuracy equals macro-averaged class recall
on the balanced samples. Three-seed summaries are arithmetic means of
metrics evaluated separately for each checkpoint, not predictions from an
ensemble. Seed variation is reported as the sample standard deviation and
is distinguished from finite-test-sample uncertainty. All accuracy
differences are expressed in percentage points.

For signal class $s$, the rejection score is the signal-minus-QCD logit,
$d_s=z_s-z_{\mathrm{QCD}}$, which has the same ordering as
$p_s/(p_s+p_{\mathrm{QCD}})$. At target signal efficiency $\epsilon_s$,
the threshold is the $\lceil\epsilon_s N_s\rceil$-th largest signal score;
events with scores equal to the threshold pass. We record the achieved
signal efficiency and QCD passing count, and define rejection as
$R_s=N_{\mathrm{QCD}}/N_{\mathrm{QCD,pass}}$ when the denominator is
nonzero. These are operating points measured on the evaluated split, not
deployment thresholds calibrated on validation. The main operating points
are 30\% and 50\% signal efficiency. The full-test evaluation also records
75\%, 99\%, and 99.5\%, with the same conventions for every model.

Rejection is calculated over the complete evaluation sample for each seed,
not averaged over file-level operating points. Percentage improvements use
the ratio of the three-seed mean rejection to the corresponding baseline
mean. If any seed passes no QCD events, we report the count and saturation
rather than a finite three-seed mean or a claim of infinite population
rejection. For the full test, 95\% Wilson intervals quantify background
count uncertainty conditional on the measured signal threshold; they do
not include uncertainty in that threshold and are not full ROC confidence
intervals.

The full-test protocol additionally computes paired accuracy intervals
against the same-seed baseline, and for selected EdgeValue against the
selected layerwise model without EdgeValue. A class-stratified bootstrap
resamples paired jet identities with class counts fixed. We use 10,000
replicates, a fixed PCG64 random seed shared across comparisons, and
linearly interpolated 2.5\% and 97.5\% quantiles. At full-test scale, multinomial sampling of per-class paired
loss/tie/win counts yields the same bootstrap distribution as resampling
event identities without materializing the index arrays. These intervals
condition on the trained checkpoints; seed standard deviations separately
describe the observed training variation.

Secondary diagnostics include per-class recall, one-versus-rest AUC,
confusion matrices, multiclass Brier score, and 15-bin top-label expected
calibration error. Calibration bins are equally spaced on $[0,1]$, left
inclusive and right exclusive except for the final bin, which includes 1;
empty bins contribute zero. The error is the occupancy-weighted absolute
difference between mean confidence and accuracy in each bin. Checkpoints,
normalizers, split identities, and source versions are recorded in the
evaluation manifests so that the original 500,000-jet and official
20-million-jet evaluations remain distinguishable.
